% !PRINTED 187
\weberpage{187}
%<<CONT
\weq{1}{I(a') \doteq r^{\lambda} I(a)}
holds. The exponent $\lambda$ we shall call the \gesperrt{weight} of the
invariant. If this exponent is zero, then $I$ is called an
\gesperrt{absolute invariant}.

One mostly considers only \gesperrt{integral rational invariants},
that is, functions $I$ which depend rationally on the coefficients $a$
and are moreover integral homogeneous functions of
them. The simultaneous invariants are required to be homogeneous
with respect to the coefficients of each of the forms on which
they depend.

Besides the invariants there are also functions which, in addition to
the coefficients, still contain the variables themselves, but in other
respects have the same properties as the invariants (1), ``the invariant
property''; thus, if $C(x, a)$ is such a
function,

\weq{2}{C(x', a') \doteq r^{\lambda} C(x, a).}

Such functions are called \gesperrt{covariants of the form $F(x)$}
% [sic] The print omits the sentence-ending period after F(x).
The functional determinant is a simultaneous covariant of a
system of $m$ functions, and the Hesse determinant
of a function of higher than the second degree is a
covariant of a single form. Among the covariants too
one considers preferably integral rational and homogeneous
functions of $x$.

If $m$ is the number of the variables $x$, $n$ the degree of the function
$F(x)$, and $\nu$ and $\mu$ the degrees of $C(x, a)$ with respect to the $x$ and
$a$, then among the numbers $\lambda, \mu, \nu, m, n$ there holds the relation

\weq{3}{n\mu = m\lambda + \nu,}

which one proves by comparing the degree of the right and
left sides of (2) with respect to the substitution coefficients.
The $a'$ are namely of degree $n$ in the substitution coefficients,
$r$ of degree $m$, $x$ of the first degree, from which
(3) follows. The corresponding formula for the invariants
is obtained by setting $\nu = 0$; indeed, the
invariants may altogether be regarded as a special case of the
covariants.

A general law for the formation of invariants and
covariants, which we have already used in the particular cases of \S.\,59,
we shall now briefly discuss.
